\scanpage{217}
\mspage{216}
Lien avec la variété des modules pour des
motifs polarisés rigidifiés : si \deleted{\unsure{le}} $M$ est
une telle variété modulaire, alors
un morphisme $X\mathrel{\text{---}}M$ est déterminé
quand on connaît sa valeur \uline{en un point},
plus l'application $\pi_1(X)\mathrel{\text{---}}\pi_1(M)$ \ldots

C'est une propriété intrinsèque importante de
ces $M$, et une propriété de cette
nature mérite d'être étudiée pour ses
propres mérites \ldots Peut-être peuvent elles
\unsure{ttes} être \deleted{caractérisées}\inserted{interprétées} comme des \deleted{schémas}\inserted{variétés}
modulaires pour certains types de
motifs.

\figure{\begin{tikzpicture}\draw (0,0)--(3,0);\end{tikzpicture}}

Introduire la notion de schéma \inserted{\unsure{motiviquement}} pointé
\inserted{\figure{\begin{tikzpicture}\node[draw,ellipse,text width=3cm,align=center] {(à val. dans un anneau $h$)};\end{tikzpicture}}}, et
morphismes de tels, qui définissent des
morphismes correspondants pour les
$G(X,\xi)\mathrel{\text{---}}G(Y,\eta)$, donc pour
\diagram{\begin{tikzcd}[column sep=large]
G(X,\xi)(h)=\pi_1(X,\xi)\arrow[r,no head]&G(Y,\eta)(h)=\pi_1(Y,\eta).
\end{tikzcd}}
